% Generated from Review/corrected_text.md (the reviewed source text).

\bakhsh{د. روند برگزاری و امتیازدهی بازی‌ها}

\band{۱. محدودیت زمانی}

برای اندازه‌گیری زمان تمام بازی‌های مسابقات از ساعت‌های دیجیتال استفاده خواهد شد. اگر ساعت دیجیتال به تعداد کافی در دسترس نباشد، از ساعت‌های آنالوگ نیز استفاده می‌شود. در دو روز اول، هر بازیکن برای انجام تمام حرکت‌های خود در هر بازی ۳۰ دقیقه زمان خواهد داشت. در بازی‌های پلی‌آف پایان این دو روز نیز هر بازیکن در هر بازی ۳۰ دقیقه زمان خواهد داشت. در تمام بازی‌های نیمه‌نهایی، هر بازیکن در هر بازی ۳۵ دقیقه زمان خواهد داشت. در فینال‌ها هر بازیکن در هر بازی ۴۰ دقیقه زمان خواهد داشت. در تمام پلی‌آف‌های مقام سوم، هر بازیکن ۳۵ دقیقه زمان خواهد داشت. اگر به هر دلیلی نتوان این محدودیت‌های زمانی را اعمال کرد، بازی‌ها هرگز نباید با زمانی کمتر از ۲۵ دقیقه برای هر بازیکن در هر بازی برگزار شوند.

\band{۲. داور میز}

هر زمان که امکان‌پذیر باشد، رئیس کل مسابقات باید یکی از مسئولان برگزاری را به عنوان داور میز برای یکی از بازی‌های مسابقات تعیین کند. داور میز مسئول ثبت حرکات بازی، پخش زنده بازی در صورت نیاز، نظارت بر رعایت قوانین اتللو و اطمینان از روند منظم و بی‌حاشیه بازی خواهد بود. پیش از آغاز بازی، داور میز باید مطمئن شود که هر دو بازیکن برای تکمیل بازی به اندازه کافی مهره دارند. تنظیم ساعت و قرار دادن آن در سمت راست بازیکنِ دارای مهره‌های سیاه (مگر آنکه هر دو بازیکن بر خلاف آن توافق کنند) وظیفه داور است. در طول بازی، داور میز باید تلاش کند حرکات بازی را بدون آنکه نیاز باشد بازیکنان حرکت خود را اعلام کنند، ثبت کند. هرگز نباید اجازه داد که ثبت حرکات بازی در تمرکز بازیکنان اختلال ایجاد کند. اگر در طول بازی مهره‌های یکی از بازیکنان تمام شود، داور میز باید ترتیبی دهد که تعداد کافی مهره از حریف به او منتقل شود. داور برای این کار مجاز است ساعت را به مدت لازم متوقف کند.

اگر برای یکی از بازی‌های مسابقات داور میز تعیین نشده باشد، هر دو بازیکن موظف به ثبت حرکات بازی هستند. در صورتی که کمک دیگری لازم باشد، هر یک از بازیکنان می‌تواند ساعت را متوقف کند و رئیس کل مسابقات یا یکی از دستیاران تعیین‌شده را فرا بخواند. اگر اختلافی پیش آید، رئیس کل مسابقات می‌تواند برای تصمیم‌گیری درباره اختلاف، برگه ثبت حرکات را از هر دو بازیکن بخواهد. در این هنگام، اگر بازیکنی حرکات بازی را ثبت نکرده باشد، می‌تواند بلافاصله آن را از حافظه بازسازی کند. اگر فقط یکی از بازیکنان بتواند برگه ثبت حرکاتی ارائه دهد که به وضعیت فعلی صفحه منجر شود، از همان برگه برای حل اختلاف استفاده خواهد شد.

\band{۳. اشیای روی صفحه بازی}

هیچ شیئی نباید به هیچ شکلی روی صفحه بازی قرار گیرد، به استثنای مهره‌های اتللو که وضعیت صفحه را تشکیل می‌دهند و یک مهره اتللو که بازیکن برای انجام حرکت خود از آن استفاده می‌کند.

\band{۴. انجام حرکت}

بازیکن اجازه ندارد در زمان حریفش حرکت خود را انجام دهد. اگر بازیکنی چنین کند، حریفش می‌تواند رئیس کل مسابقات را فرا بخواند تا هر اقدام لازم را انجام دهد (برای مثال افزایش یا کاهش زمان بازیکنان). اگر بازیکنی فراموش کند پس از حرکتش ساعت خود را بزند، حریفش مجاز است و حتی تشویق می‌شود که به او یادآوری کند. هیچ تماشاگری، به جز داور میز، رئیس کل مسابقات یا یکی از دستیارانش، حق ندارد زدن ساعت را به بازیکن یادآوری کند. اگر چنین اتفاقی بیفتد، رئیس کل مسابقات باید آن تماشاگر را از سالن مسابقات اخراج کند و می‌تواند مطابق بند هـ.۲ جریمه‌ای برای او در نظر بگیرد.

بازیکن اگر با قصد آشکار انجام حرکت در خانه‌ای، صفحه را با دست یا با مهره لمس کند، متعهد به انجام حرکت در آن خانه می‌شود. در این حالت، به شرط آنکه بازیکن در آن خانه حرکت مجاز داشته باشد، باید در اولین خانه‌ای که لمس کرده است حرکت کند و نمی‌تواند نظر خود را عوض کند و در خانه دیگری بازی کند، حتی اگر حریفش به او اجازه پس گرفتن حرکت را بدهد. هدف از این قانون، جلوگیری از پرت شدن حواس حریف یا مسدود شدن دید او به صفحه بازی توسط بازیکنی است که مردد است. اگر داور میز تشخیص دهد که بازیکن خانه را ناخواسته لمس کرده است، بازیکن ملزم به رعایت این قانون نیست.

به محض آنکه بازیکن متعهد به انجام حرکتی شد، باید تمام مهره‌های گرفته‌شده را برگرداند. او باید این کار را فقط با یک دست انجام دهد و سپس با همان دست ساعت خود را بزند (حتی اگر حریفش پس از آن پاس دهد). اگر بازیکن پس از انجام حرکت فراموش کند ساعت خود را بزند، داور میز باید به او یادآوری کند. بازیکن باید با زدن ساعت خود اعلام کند که پاس می‌دهد.

\band{۵. حرکات غیرمجاز}

\zirband{\LRE{۵\kern0pt.\kern0pt۱} روند رسیدگی}

در صورت بروز اختلافی که با قوانین زیر قابل حل نباشد (برای مثال اختلاف نظر درباره وضعیت صفحه، اختلاف نظر درباره اینکه آیا مهره‌ای به اشتباه برگردانده شده است یا نه و ...) در میزی که داور میز ندارد، بازیکنان باید زمان باقی‌مانده روی ساعت‌هایشان را یادداشت کنند و بی‌درنگ رئیس کل مسابقات یا یکی از دستیارانش را برای حل اختلاف فرا بخوانند.

داور میز مسئول اجرای قوانین اتللو است. اشتباه در حرکتی که به تازگی انجام شده و هنوز به آن پاسخ داده نشده است، ممکن است از سوی هر یک از بازیکنان، رئیس کل مسابقات یا یکی از دستیارانش به اطلاع داور میز برسد (اما نه از سوی تماشاگران؛ تماشاگر باید رئیس کل مسابقات را فرا بخواند تا او درباره اقدام لازم تصمیم بگیرد). به محض اینکه داور میز (یا رئیس کل مسابقات یا یکی از دستیارانش) متوجه چنین اشتباهی شود یا از آن مطلع شود، باید به ترتیب زیر عمل کند:

\begin{bandha}
  \item[الف)] ساعت بازیکنی را که اشتباه کرده است دوباره به کار اندازد.
  \item[ب)] خودش اشتباه را اصلاح کند.
  \item[ج)] ساعت بازیکنی را که نوبت حرکت اوست دوباره به کار اندازد.
\end{bandha}

اگر میزی داور میز نداشته باشد، اشتباه در حرکتی که به تازگی انجام شده و هنوز به آن پاسخ داده نشده است، می‌تواند از سوی هر یک از بازیکنان اعلام شود و سپس روند زیر طی می‌شود:

\begin{bandha}
  \item[الف)] ساعت بازیکنی که اشتباه کرده است، دوباره به کار انداخته شود.
  \item[ب)] از آن بازیکن خواسته شود که اشتباه را اصلاح کند.
  \item[ج)] ساعت بازیکنی که نوبت حرکت اوست، دوباره به کار انداخته شود.
\end{bandha}

\zirband{\LRE{۵\kern0pt.\kern0pt۲} پذیرش وضعیت}

فقط آخرین حرکت انجام‌شده قابل اعتراض است. وقتی بازیکنی حرکت جدیدی انجام می‌دهد، به طور ضمنی وضعیت صفحه را همان‌گونه که هست پذیرفته است و نه او و نه هیچ‌یک از مسئولان برگزاری نمی‌توانند بعداً به دنبال اصلاح اشتباه قبلی باشند.

\zirband{\LRE{۵\kern0pt.\kern0pt۳} وضعیت ناخوانا}

بازیکنان باید حرکات خود را طوری انجام دهند و مهره‌ها را طوری در وسط خانه‌ها قرار دهند که وضعیت صفحه به وضوح قابل خواندن باشد. اگر به نظر داور میز وضعیت صفحه ناخوانا شود، او باید هر اقدامی را برای اصلاح آن انجام دهد (از جمله متوقف کردن ساعت‌ها و بازچینی بازی از روی برگه ثبت حرکات). او همچنین می‌تواند از رئیس کل مسابقات بخواهد که جریمه‌ای برای بازیکن خطاکار در نظر بگیرد.

\band{۶. اتمام وقت}

بازیکن بازی را می‌بازد، مگر اینکه بتواند پیش از افتادن پرچم ساعتش، حرکت نهایی خود را انجام دهد، تمام مهره‌های لازم را برگرداند و ساعت خود را متوقف کند. اگر پرچم ساعت بازیکنی بیفتد، او بازی را می‌بازد و حریفش بدون توجه به وضعیت صفحه برنده اعلام می‌شود. هر دو بازیکن دو دقیقه زمان اضافه دریافت می‌کنند. بازی سپس ادامه می‌یابد و نتیجه نهایی واقعی ملاک خواهد بود، با این تفاوت که برای بازیکنی که وقتش تمام نشده بود، دست‌کم یک برد ۳۳ بر ۳۱ تضمین شده است. اگر پرچم ساعت بازیکنی که وقتش تمام شده بود برای بار دوم بیفتد، نتیجه بازی ۶۴ بر ۰ ثبت می‌شود. اگر پرچم ساعت بازیکنی که وقتش تمام نشده بود، پس از دریافت دو دقیقه زمان اضافه بیفتد، بازیکنی که ابتدا وقتش تمام شده بود، بازی را با نتیجه ۳۱ بر ۳۳ می‌بازد.

\band{۷. ساعت معیوب}

اگر در طول بازی مشخص شود که ساعت مورد استفاده معیوب است، داور میز (یا اگر داور میز حضور ندارد، بازیکنان) رئیس کل مسابقات را فرا می‌خوانند تا هر اقدامی را که مناسب بداند انجام دهد.

\band{۸. امتیازدهی}

در پایان بازی، اگر هر دو بازیکن حرکات خود را در زمان مجاز کامل کرده باشند، برنده بازیکنی است که در پایان بازی تعداد بیشتری مهره از رنگ خود روی صفحه دارد. نتیجه رسمی بازی با شمارش مهره‌های هر رنگ روی صفحه تعیین می‌شود و خانه‌های خالی به حساب برنده گذاشته می‌شوند. در صورت تساوی، نتیجه همیشه ۳۲ بر ۳۲ خواهد بود. اگر وقت یکی از بازیکنان تمام شود، نتیجه به شکلی که در بند د.۶ آمده است تعیین می‌شود. اگر بازیکنی به هر دلیلی نتواند یا نخواهد بازی خود را تمام کند، نتیجه رسمی بازی ۶۴ بر ۰ به نفع حریف او خواهد بود. پس از پایان هر بازی، هر دو بازیکن باید برگه نتیجه را امضا کنند تا نشان دهند با نتیجه نهایی موافق‌اند و آن را به یکی از مسئولان برگزاری مسابقات تحویل دهند.

\band{۹. انتخاب رنگ یا تساوی-برد}

وقتی بازیکنی حق انتخاب میان رنگ بازی یا تساوی-برد را دارد، می‌تواند:

\begin{bandha}
  \item[الف)] خودش رنگ بازی را انتخاب کند؛ در این صورت اگر بازی مساوی شود، حریفش برنده اعلام می‌شود.
  \item[ب)] اجازه دهد حریفش رنگ را انتخاب کند؛ در این صورت اگر بازی مساوی شود، خودش برنده اعلام می‌شود.
\end{bandha}

وقتی بازیکنی می‌تواند رنگ خود را برای بازی انتخاب کند، باید دست‌کم پنج دقیقه پیش از آغاز بازی، انتخاب خود را به رئیس کل مسابقات و حریفش اعلام کند.

\band{۱۰. قوانین مربوط به بازیکنان دارای نیازهای ویژه}

اگر رئیس کل مسابقات تشخیص دهد که بازیکنی نیازهای ویژه‌ای دارد، آن بازیکن می‌تواند دستیاری داشته باشد که برخی یا همه کارهای زیر را (بسته به نیاز بازیکن) برای او انجام دهد:

\begin{itemize}
  \item زدن ساعت
  \item اطلاع دادن زمان باقی‌مانده به بازیکن
  \item اطلاع دادن به بازیکن هنگامی که حریف نوبت خود را انجام داده است
  \item قرار دادن و برگرداندن مهره‌ها (پس از آنکه بازیکن حرکت مورد نظر خود را به روشی مورد تأیید رئیس کل مسابقات مشخص کرد)
  \item هر نیاز دیگری که رئیس کل مسابقات برای فراهم کردن فرصتی منصفانه و متوازن برای شرکت بازیکن در مسابقات لازم بداند.
\end{itemize}

دستیار به هیچ وجه اجازه ندارد به شکل دیگری در بازی دخالت کند یا به بازیکن راهنمایی دهد. اگر لازم باشد بازیکنِ دارای آسیب بینایی برای تشخیص وضعیت صفحه، صفحه را با دست لمس کند، هیچ جریمه‌ای برای این کار متوجه او نیست. با این حال، حریف می‌تواند درخواست کند که یک صفحه بازی اصلی وجود داشته باشد که همواره بدون مانع باشد؛ در این صورت، نوبت بازیکنِ دارای آسیب بینایی تنها زمانی کامل می‌شود که دستیارش حرکت را روی آن صفحه انجام دهد و سپس ساعت زده شود. در این حالت، پنج دقیقه زمان اضافه به ساعت بازیکنِ دارای آسیب بینایی افزوده می‌شود تا زمان لازم برای انتقال حرکت بین صفحه‌ها را داشته باشد.
